% year: תשע"ט
% type: מבחן
% semester: א
% moed: מיוחד
% number: 6א
% points: 10
% categories: improper-integrals, antiderivatives
% source: exams/מבחנים/תשעט/מועד מיוחד תשעט.pdf

\begin{question}
האם האינטגרל הלא אמיתי הבא $\displaystyle\int_{-\infty}^{-2}\frac{dx}{x^2-2x}$ מתכנס? אם כן, מצאו אותו.
\end{question}

\begin{hint}
פרקו לשברים חלקיים: $\frac{1}{x(x-2)}=\frac12\left(\frac{1}{x-2}-\frac1x\right)$.
\end{hint}

\begin{solution}
בקטע $(-\infty,-2]$ המכנה $x^2-2x=x(x-2)>0$ (שני הגורמים שליליים), ולכן הפונקציה רציפה וחיובית שם, והבעיה היחידה היא הגבול $-\infty$. לפי ההגדרה,
\[ \int_{-\infty}^{-2}\frac{dx}{x^2-2x} = \lim_{R\to-\infty}\int_R^{-2}\frac{dx}{x(x-2)}. \]
פירוק לשברים חלקיים: $\frac{1}{x(x-2)}=\frac{A}{x}+\frac{B}{x-2}$, ומ-$1=A(x-2)+Bx$ נקבל ($x=0$) $A=-\frac12$ ו-($x=2$) $B=\frac12$. לכן
\[ \int\frac{dx}{x(x-2)} = \frac12\ln|x-2|-\frac12\ln|x| + C = \frac12\ln\left|\frac{x-2}{x}\right| + C. \]
לפיכך
\[ \int_R^{-2}\frac{dx}{x(x-2)} = \frac12\ln\frac{-4}{-2} - \frac12\ln\left|\frac{R-2}{R}\right| = \frac12\ln2-\frac12\ln\left|1-\frac2R\right|. \]
כאשר $R\to-\infty$, $1-\frac2R\to1$ ומרציפות הלוגריתם $\ln\left|1-\frac2R\right|\to0$. לכן הגבול קיים וסופי:
\[ \int_{-\infty}^{-2}\frac{dx}{x^2-2x} = \frac{\ln2}{2}. \]
\textbf{תשובה:} האינטגרל מתכנס וערכו $\frac12\ln2$.
(אפשר גם לראות את ההתכנסות ממבחן ההשוואה הגבולי עם $\frac1{x^2}$, אך החישוב הישיר נותן גם את הערך.)
\end{solution}
